% The cover letter as ACS wants it: one PDF.
%
% SOURCE OF RECORD for the letter that is uploaded.  paper/cover_letter.md is the working
% document -- it also carries the text for JCIM's own submission form, the pre-submission
% checklist, and the reasoning behind the reviewer list, none of which belongs in a letter to an
% editor.  What is here is Part 1 of that file plus the reviewer names, and the two must be kept
% in step by hand: if you edit the letter in one, edit it in the other.
%
%   xelatex cover_letter.tex
%
\documentclass[11pt]{article}
\usepackage[a4paper,margin=25mm]{geometry}
\usepackage{fontspec}
% NOT STIX HERE, and for a reason worth writing down.  The manuscript is set in STIX Two Text,
% but on this machine that family ships as one variable font whose Bold is a named instance
% XeTeX will not select, and whose Italic lives in a second file under the same family name.
% \setmainfont{STIX Two Text} therefore compiles happily and silently renders every \textbf at
% regular weight -- the manuscript title and the reviewer names came out unbolded, and pdffonts
% confirmed no bold face was embedded at all.  Naming the faces explicitly fails outright,
% because "STIX Two Text Bold" is not a font name.  Times New Roman resolves all four faces,
% which pdffonts confirms, and a cover letter in Times is unremarkable.
\setmainfont{Times New Roman}
\usepackage{unicode-math}
\setmathfont{STIX Two Math}   % text is Times; the maths would otherwise be Computer Modern
\usepackage{microtype}
\usepackage[hidelinks]{hyperref}
\usepackage{parskip}
\usepackage{enumitem}
\setlist[itemize]{leftmargin=1.2em,itemsep=2pt,topsep=3pt}
\pagestyle{empty}

\begin{document}

\noindent
{\large Nikita L. Polomoshnov}\\[2pt]
Faculty of Bioengineering and Bioinformatics, Lomonosov Moscow State University\\
V.~N. Orekhovich Institute of Biomedical Chemistry, Moscow, Russia\\
\href{mailto:nikitapol@fbb.msu.ru}{nikitapol@fbb.msu.ru}

\vspace{4mm}
\noindent 3 September 2026

\vspace{4mm}
\noindent
To the Editors\\
\emph{Journal of Chemical Information and Modeling}

\vspace{4mm}
\noindent Dear Editors,

We submit for consideration as an \textbf{Article} the manuscript \textbf{``Ground Truth That Does
Not Ground: Substituting Reference Solvent $\sigma$-Profiles Degrades COSMO-SAC Solubility
Prediction''} by N.~L. Polomoshnov and A.~V. Rudik.

Physics-informed models are built so that a prediction passes through a quantity chemistry already
tabulates, and the design carries an assumption that is rarely tested: that handing the model the
tabulated value, in place of the one it learned, should improve it. We test that assumption
directly. A graph neural network predicts the $\sigma$-profiles a fixed COSMO-SAC layer consumes,
and quantum chemistry tabulates the same profiles independently, so the substitution can be made
without retraining anything.

It makes the model worse. Over five seeds the mean absolute error on solid--liquid solubility rises
from $1.93$ to $2.34$ in $\ln x_2$, about $21\%$, with no overlap between the per-seed values of the
two arms, and the mean per-solute Spearman correlation over solvents falls by $0.185$, a loss in the
ranking that solvent selection actually uses. Supervising the same network against the same
database during training, the other operation the word ``grounding'' covers, gives gains that
straddle zero, and we report it as having no established sign rather than as a benefit. The
operation that fails is the one available in practice, because it is the one that needs no
retraining.

The constructive half of the paper is a way to tell in advance where this will happen. Using
infinite-dilution activity coefficients we separate a deployed thermodynamic model's own error from
what its inputs leave unresolved, search fifty-nine chemical strata, and find the model's own error
dominant on one solvent chemistry. We state plainly which of those bounds are established and which
are the instrument's output on a set too small to certify, including one that an independent set
reproduces at half its size without meeting our own stability rule.

We think this belongs in JCIM for two reasons. It concerns how a widely used cheminformatics
construction behaves when a learned component meets a fixed physical map, which is a modelling
question rather than a chemistry-specific one; and it is written to the standard of the Journal's
editorial on method and data sharing (\emph{J. Chem. Inf. Model.} \textbf{2015}, \emph{55}, 719).
Every number in the manuscript and the Supporting Information is produced by a deposited script
from a deposited artifact; the code, checkpoints, processed splits and per-arm predictions are
released; and the two decomposition sets are deposited as row-level machine-readable tables so a
reader can recompute the central claim without running our code. Twelve defects and limits of the
work are disclosed by us, each where its numbers print, and none of them is reported by anyone
else.

This manuscript has not been published elsewhere and is not under consideration by any other
journal. All authors have read and approved the final version of the manuscript. We declare no
conflicts of interest. This study was supported by the Russian Science Foundation, project
No.~25-25-00148.

\vspace{2mm}
\noindent
We would be glad to answer any questions.

\vspace{4mm}
\noindent
Sincerely,\\[2pt]
Nikita L. Polomoshnov, on behalf of both authors

\vspace{6mm}
\noindent\textbf{Suggested reviewers}

\begin{itemize}
\item \textbf{Fabian Jirasek}, Laboratory of Engineering Thermodynamics, RPTU
      Kaiserslautern-Landau, Germany. \href{mailto:fabian.jirasek@rptu.de}{fabian.jirasek@rptu.de}
      \\ Machine-learning models for activity coefficients under thermodynamic constraints.
\item \textbf{André Bardow}, Energy and Process Systems Engineering, ETH Zürich, Switzerland.
      \href{mailto:abardow@ethz.ch}{abardow@ethz.ch}
      \\ Learned predictors for limiting activity coefficients, and physics-guided models
      constrained to thermodynamic consistency.
\item \textbf{Ian H. Bell}, Applied Chemicals and Materials Division, NIST, Boulder, Colorado, USA.
      \href{mailto:ian.bell@nist.gov}{ian.bell@nist.gov}
      \\ Author of the open-source reference COSMO-SAC implementation this work validates against.
\item \textbf{Chieh-Ming Hsieh}, Department of Chemical and Materials Engineering, National Central
      University, Taoyuan City, Taiwan. \href{mailto:hsiehcm@ncu.edu.tw}{hsiehcm@ncu.edu.tw}
      \\ Author of COSMO-SAC-2010 and of the dispersion term the manuscript implements.
\item \textbf{William H. Green}, Department of Chemical Engineering, Massachusetts Institute of
      Technology, USA. \href{mailto:whgreen@mit.edu}{whgreen@mit.edu}
      \\ SolProp and FastSolv, the external points of comparison used here.
\item \textbf{Hans Hasse}, Laboratory of Engineering Thermodynamics, RPTU Kaiserslautern-Landau,
      Germany. \href{mailto:hans.hasse@rptu.de}{hans.hasse@rptu.de}
      \\ Engineering thermodynamics of mixtures, and hard-constrained models for activity
      coefficients.
\item \textbf{Kai Sundmacher}, Otto-von-Guericke University Magdeburg and Max Planck Institute for
      Dynamics of Complex Technical Systems, Germany.
      \href{mailto:sundmacher@mpi-magdeburg.mpg.de}{sundmacher@mpi-magdeburg.mpg.de}
      \\ Graph networks for the temperature dependence of activity coefficients.
\end{itemize}

\end{document}
